\chapter{RESUMO}\label{chap:resumo}

Este documento apresenta a estrutura básica do texto da Proposta de Trabalho Final do Curso de Ciência da Computação e Trabalho de Conclusão do Curso de Engenharia de Computação da Universidade de Passo Fundo. A elaboração da PTC tem elementos textuais fixos e obrigatórios. Para detalhes sobre o conteúdo a ser desenvolvido em cada item da proposta, converse com o professor da Disciplina de Projeto Aplicado ou com seu futuro orientador. \ul{Por favor, leia com atenção todas as instruções que estão descritas ao longo deste documento, bem como os comentários dentro do projeto \LaTeX.} Um resumo consistente de um texto científico deve conter todos os elementos abordados no texto. Procure não repetir frases do texto no resumo (e também no texto). Evite o uso de siglas no resumo, mas se forem imprescindíveis, coloque sua descrição por extenso antes. Para uma Proposta de Trabalho de Conclusão ou Trabalho Final recomenda-se falar da motivação, dos objetivos, dos achados da revisão de literatura e da metodologia a ser desenvolvida no resumo. O \textbf{Resumo} deve ser um texto (parágrafo) único, sem quebra de linhas.

